%!TeX TxS-program:compile =txs:///pythontex
\documentclass{mwart}
\usepackage[utf8]{inputenc}
\usepackage{polski}
\usepackage{pythontex}
\usepackage{graphicx}
\usepackage{pgf}
\usepackage{lmodern,microtype}
\usepackage{amsfonts}
\usepackage{float}

\addtolength{\topmargin}{-1cm}
\addtolength{\oddsidemargin}{-2cm}
\addtolength{\textwidth}{4cm}
\addtolength{\textheight}{0cm}

\def \R{\mathbb{R}}

\author{Marta Hałas 282264}
\title{\textbf{Lista 2}}

\begin{document}
	\maketitle
	\medskip
	
	\section{\textbf{Zadanie 1}}\label{Zadanie 1}
	
	\begin{pysub}
		
		Skoczek pustynny Alfred w tym stanie nie trafi do swojego domu, ani też nie będzie nocował pod gołym niebem. Będzie on nocował pod dachem swojego kolegi, aby tam dotrzeć musi wykonać po wyjściu z baru 3 kroki w przód o długości 2.28m i 3 kroki wstecz o długości 84cm.
		
	\end{pysub}
	
	
	\section{\textbf{Zadanie 2}}\label{Zadanie 2}

	
	\begin{figure}[H]
		\begin{center}
			\includegraphics[width=15cm]{WYKRES2LISTASTOP.png}
			\caption{Wykres: Czas działania implementacji dla różnych wartości parametrów ${n\choose k}$}
		\end{center}
	\end{figure}

	\medskip

	\begin{pysub}
	
		Najszybszą implementacją okazała się Newton silnia, drugie miejsce zajęła Newton rekurencja. Natomiast najwolniejsza okazała się Newton iteracja.
		 
	\end{pysub}
	
	
	\begin{figure}[H]
		\begin{center}
			\includegraphics[width=15cm]{2opcja2wykees.png}
			\caption{Wykres: Czas działania implementacji dla różnych wartości parametrów ${n\choose k}$}
		\end{center}
	\end{figure}
	
	\medskip
	
	\begin{pysub}
		
		Najszybszą implementacją okazała się Newton silnia, drugie miejsce zajęła Newton iteracja, chociaż dla n mniejszych od 9 była najwolniejsza. Newton rekurencja dla n większego od 9 okazała się najwolniejsza. Dla n mniejszych od 7 Newton reukrencja konkurowała z implementacją Newton silnia. Dla n=7, n=8, n=9 drugą najszybszą implementacją była implementacja Newton rekurencja.
		
	\end{pysub}
	
\end{document}